\section{The Interface Diagnostic}
\label{sec:method}
% ============================================================================

\subsection{Problem statement}

\begin{problem}[Interface diagnosis]
\label{prob:main}
\textbf{Input:} a code $\Phi\in\R^{d\times F}$ with $F>d$ (or a trained network from which an
effective code is read off); a readout family $\mathcal G$; a sparse-state model, here
uniformly random supports of size $s$ drawn from a grid $\mathcal S$; a threshold $\theta$; a
trial budget $T$; a confidence level $1-\alpha$.

\textbf{Output:} (i) the attainment ratio $r=\widehat W/W(F,d)\ge1$ and the maximum-form ratio
$r_{\max}$, together with its factorisation $R_{\mathrm{geom}}\times R_{\mathrm{readout}}$
against the code's own floor (Section~\ref{sec:codefloor}); (ii) for each $s\in\mathcal S$, the
empirical $p_{\mathrm{rec}}(s)$ with a Wilson interval and the exact per-coordinate RMS error
$\rho_{\mathrm{lin}}(s)$ of the calibrated linear interface; (iii) the criterion-separation
witness $(s_{95},\rho_{\mathrm{lin}}(s_{95}))$; (iv) for each feature in a chosen subset, the
collision frontier $\kappa_i$ and hence the largest sparsity at which any affine probe can
detect that feature (Algorithm~\ref{alg:frontier}).
\end{problem}

The output is a diagnosis, not a score to be optimised, and the four parts answer different
questions. A code with $r$ close to $1$ is Welch-optimal in the mean-square sense, but only
$R_{\mathrm{readout}}$ says anything about the decoder; a large $s_{95}$ with
$\rho_{\mathrm{lin}}(s_{95})$ well above $d^{-1/2}$ says the code supports support recovery in a
regime where analog reconstruction from the same map is useless; and $\kappa_i$ is the only part
that is per-feature and exact rather than aggregate and empirical, which is what makes it usable
as a statement about an individual feature rather than about the code on average.

\subsection{Algorithm 1: floor attainment}

\begin{algorithm}[t]
\caption{\textsc{FloorAttainment} --- how close the readout interface sits to the floor}
\label{alg:one}
\begin{algorithmic}[1]
\Require code $\Phi\in\R^{d\times F}$, $F>d$; readout $G\in\R^{F\times d}$ or a rule in
  $\{\textsc{tied},\textsc{pinv},\textsc{lsq}\}$
\Ensure attainment ratio $r$ (and $r_{\max}$ if requested), or \textsc{Degenerate}
\State $G\gets\Phi^{\top}$, $\Phi^{+}$, or the least-squares fit, per the rule
\State $D_{ii}\gets\langle g_i,\phi_i\rangle$ for $i=1,\dots,F$
  \Comment{the diagonal of $M=G\Phi$, in $O(Fd)$; $M$ itself is not formed}
\If{$\min_i|D_{ii}|<\epsilon_{\mathrm{gain}}$} \Return \textsc{Degenerate}
  \Comment{some feature is read out with vanishing gain}
\EndIf
\State $\bar G\gets D^{-1}G$ \Comment{gain normalisation; see Corollary~\ref{cor:calibfree}}
\State $H\gets\bar G^{\top}\bar G$;\quad $\Sigma\gets\Phi\Phi^{\top}$
  \Comment{both $d\times d$; $O(Fd^2)$}
\State $\widehat W\gets\dfrac{\operatorname{tr}(H\Sigma)-F}{F(F-1)}$
  \Comment{exact, by Lemma~\ref{lem:suffstat}}
\State $W\gets(F-d)/(d(F-1))$ \Comment{Theorem~\ref{thm:floor}}
\If{the maximum form is requested}
  \State form $\widetilde M\gets\bar G\Phi$ and set
    $\widehat\mu\gets\max_{i\neq j}|\widetilde M_{ij}|$
    \Comment{$O(F^2d)$ time, $O(F^2)$ memory}
\EndIf
\State \Return $r\gets\widehat W/W$ \ (and $r_{\max}\gets\widehat\mu/\sqrt{W}$ if computed)
\end{algorithmic}
\end{algorithm}

Step~3 is not a numerical guard but a semantic one: without it, an interface can be made to
look arbitrarily good by shrinking $G$, which is exactly the confusion that
Remark~\ref{rem:calib} dispels. By Theorem~\ref{thm:floor}, a correct implementation always
returns $r\ge1$; we use that as a runtime assertion, and no violation was observed in
\EoneMeasurements{} gain-normalised interfaces (Section~\ref{sec:res-e1}).

The mean-square statistic that the theorem actually constrains never requires the interface to
be built. The maximum statistic does, and we say so rather than folding it into the same cost.

\begin{lemma}[Mean-square cross-talk from $d\times d$ statistics]
\label{lem:suffstat}
Let $\bar G=D^{-1}G$ with $D_{ii}=\langle g_i,\phi_i\rangle\neq0$, let
$\widetilde M=\bar G\Phi$, and set $H=\bar G^{\top}\bar G$ and $\Sigma=\Phi\Phi^{\top}$, both
in $\R^{d\times d}$. Then
\[
  \sum_{i\neq j}\widetilde M_{ij}^{2}
  \;=\;\norm{\bar G\Phi}_F^{2}-F
  \;=\;\operatorname{tr}(H\Sigma)-F .
\]
\end{lemma}

\begin{proof}
$\widetilde M$ has unit diagonal by construction, so the $F$ diagonal entries contribute
exactly $F$ to $\norm{\widetilde M}_F^2$ and the remainder is the off-diagonal sum. Cyclicity
of the trace gives $\norm{\bar G\Phi}_F^2=\operatorname{tr}(\Phi^{\top}\bar G^{\top}\bar
G\Phi)=\operatorname{tr}(\bar G^{\top}\bar G\,\Phi\Phi^{\top})=\operatorname{tr}(H\Sigma)$.
\end{proof}

Forming $H$ and $\Sigma$ costs $O(Fd^2)$ time and $O(d^2)$ memory, against $O(F^2d)$ and
$O(F^2)$ for the direct route --- a reduction of a factor $F/d$ in time and $F^2/d^2$ in
memory, with no approximation whatever. The identity is verified against the dense computation
to relative accuracy $10^{-10}$ in the accompanying tests.

\subsection{Algorithm 2: separation profile of a given code}

Algorithm~\ref{alg:two} is the diagnostic proper: it takes a code and a readout that already
exist --- a designed code, an optimised one, or the effective code of a trained network --- and
holds them fixed while the support varies. This is the only version that says anything about a
particular object, and it is the one applied to trained networks in
Section~\ref{sec:res-e5}.

\begin{algorithm}[t]
\caption{\textsc{SeparationProfile} --- support recovery versus irreducible analog error, on a
fixed code}
\label{alg:two}
\begin{algorithmic}[1]
\Require code $\Phi\in\R^{d\times F}$ and readout $G\in\R^{F\times d}$; sparsity grid
  $\mathcal S$; trials $T$; threshold $\theta$; seed $\sigma$
\Ensure $\{(s,p_{\mathrm{rec}}(s),\text{CI},\rho_{\mathrm{lin}}(s))\}_{s\in\mathcal S}$ and the
  separation witness $(s_{95},\rho_{\mathrm{lin}}(s_{95}))$
\State $s_{\max}\gets\max\mathcal S$;\quad $\bar G\gets D^{-1}G$ as in
  Algorithm~\ref{alg:one}
\State $\norm{A}_F^2\gets\operatorname{tr}(H\Sigma)-F$;\quad
       $\norm{A\mathbf 1}^2\gets\norm{\bar G\Phi\mathbf 1}^2
       -2\,\mathbf 1^{\top}\bar G\Phi\mathbf 1+F$
  \Comment{Lemma~\ref{lem:moments}; $O(Fd^2)$}
\State $e_s\gets\frac1F\!\left[(p-q)\norm{A}_F^2+q\norm{A\mathbf 1}^2\right]$,
  $p=\tfrac sF$, $q=\tfrac{s(s-1)}{F(F-1)}$, for each $s\in\mathcal S$
  \Comment{exact; Prop.~\ref{prop:uniform-energy}}
\For{$t=1$ \textbf{to} $T$}
  \State draw nested supports $S_1\subset\cdots\subset S_{s_{\max}}$ from seed $(\sigma,t)$
  \State $Z\gets\bar G\,[\,x_1,\dots,x_{s_{\max}}\,]$, $x_s=\sum_{j\in S_s}\phi_j$
    \Comment{$O(Fd\,s_{\max})$}
  \State $k_s\mathrel{+}=1$ for each $s$ with $\mathbf 1\{Z_{\cdot s}\ge\theta\}=\mathbf 1_{S_s}$
\EndFor
\State $p_{\mathrm{rec}}(s)\gets k_s/T$ with Wilson interval;\quad
       $\rho_{\mathrm{lin}}(s)\gets\sqrt{e_s}$
\State $s_{95}\gets\max\{s:\ p_{\mathrm{rec}}(s')\ge0.95\ \forall s'\le s\}$
\State \Return the profile and $(s_{95},\rho_{\mathrm{lin}}(s_{95}))$
\end{algorithmic}
\end{algorithm}

The support-recovery side of Algorithm~\ref{alg:two} is Monte Carlo over supports; the analog
side is \emph{exact} for the given code, computed once before the loop rather than estimated
inside it. Total cost is $O(Fd^2+TFd\,s_{\max})$ time and $O(Fd+d^2)$ memory: the $F\times F$
interface is never formed.

\subsection{A scaling experiment on the random-code ensemble}
\label{sec:alg-ensemble}

Algorithm~\ref{alg:ensemble} is a different object and we keep it named apart. It redraws a
fresh code at every trial, so what it characterises is the $(d,F)$ random-code \emph{ensemble},
not any individual code; it is the vehicle for the width-scaling study of
Sections~\ref{sec:res-e1} and~\ref{sec:res-e6}, and it is not a diagnosis of anything. Because
the code is regenerated block by block from a counted seed and never stored, it reaches widths
at which the code would not fit in memory.

\begin{algorithm}[t]
\caption{\textsc{EnsembleScaling} --- the same two criteria, averaged over random codes}
\label{alg:ensemble}
\begin{algorithmic}[1]
\Require $d$, $F$, sparsity grid $\mathcal S$, trials $T$, threshold $\theta$, seed $\sigma$,
  block size $B$
\Ensure the ensemble-averaged profile and $(s_{95},\rho_{\mathrm{lin}}(s_{95}))$
\State $s_{\max}\gets\max\mathcal S$;\quad $k_s\gets0$ for all $s$;\quad $e_s\gets0$ for all $s$
\For{$t=1$ \textbf{to} $T$}
  \State draw nested supports $S_1\subset\cdots\subset S_{s_{\max}}$ from seed $(\sigma,t)$
  \State $X\gets$ partial sums $[\,x_1,\dots,x_{s_{\max}}\,]$, $x_s=\sum_{j\in S_s}\phi_j$
    \Comment{$d\times s_{\max}$}
  \State $\Sigma\gets 0_{d\times d}$,\ $v\gets 0_d$
  \For{each block $\beta$ of $B$ columns of $\Phi$}
    \State regenerate $\Phi_\beta$ from seed $(\sigma,t,\beta)$ \Comment{$\Phi$ is never stored}
    \State $Z\gets\Phi_\beta^{\top}X$; mark $s$ as failed if
      $\mathbf 1\{Z_{\cdot s}\ge\theta\}$ disagrees with $\mathbf 1_{S_s}$ on this block
    \State $\Sigma\mathrel{+}=\Phi_\beta\Phi_\beta^{\top}$;\quad
           $v\mathrel{+}=\Phi_\beta\mathbf 1$
  \EndFor
  \State $k_s\mathrel{+}=1$ for each surviving $s$
  \State $\norm{A}_F^2\gets\norm{\Sigma}_F^2-F$;\quad
         $\norm{A\mathbf 1}^2\gets v^{\top}\Sigma v-2\norm{v}^2+F$
    \Comment{Lemma~\ref{lem:moments}}
  \State $e_s\mathrel{+}=\frac1F\!\left[(p-q)\norm{A}_F^2+q\norm{A\mathbf 1}^2\right]$,
    $p=\tfrac sF$, $q=\tfrac{s(s-1)}{F(F-1)}$ \Comment{Prop.~\ref{prop:uniform-energy}}
\EndFor
\State $p_{\mathrm{rec}}(s)\gets k_s/T$ with Wilson interval;\quad
       $\rho_{\mathrm{lin}}(s)\gets\sqrt{e_s/T}$
\State $s_{95}\gets\max\{s:\ p_{\mathrm{rec}}(s')\ge0.95\ \forall s'\le s\}$
\State \Return the profile and $(s_{95},\rho_{\mathrm{lin}}(s_{95}))$
\end{algorithmic}
\end{algorithm}

Both algorithms share the same asymmetry: the support-recovery side is estimated, the analog
side is computed. That is what makes the comparison a witness rather than two noisy curves ---
the analog error carries no Monte Carlo uncertainty, so any gap at $s_{95}$ is attributable to
the recovery side alone.

\begin{lemma}[Sufficient statistics of the tied interface]
\label{lem:moments}
Let $\Phi$ have unit-norm columns, $M=\Phi^{\top}\Phi$, $A=M-I_F$,
$\Sigma=\Phi\Phi^{\top}\in\R^{d\times d}$ and $v=\Phi\mathbf 1\in\R^d$. Then
\[
  \norm{A}_F^2=\norm{\Sigma}_F^2-F,
  \qquad
  \norm{A\mathbf 1}_2^2=v^{\top}\Sigma v-2\norm{v}_2^2+F .
\]
\end{lemma}

\begin{proof}
By cyclicity of the trace, $\norm{\Phi^{\top}\Phi}_F^2=\operatorname{tr}((\Phi^{\top}\Phi)^2)
=\operatorname{tr}((\Phi\Phi^{\top})^2)=\norm{\Sigma}_F^2$; the diagonal of $M$ is all ones and
contributes $F$, so subtracting $F$ leaves the off-diagonal energy, which is $\norm{A}_F^2$.
For the second identity, $A\mathbf 1=\Phi^{\top}\Phi\mathbf 1-\mathbf 1=\Phi^{\top}v-\mathbf 1$,
hence $\norm{A\mathbf 1}^2=v^{\top}\Phi\Phi^{\top}v-2\mathbf 1^{\top}\Phi^{\top}v+F
=v^{\top}\Sigma v-2\norm{v}^2+F$, using $\mathbf 1^{\top}\Phi^{\top}v=v^{\top}v$.
\end{proof}

Lemma~\ref{lem:moments} is the reason the diagnostic scales: it replaces an $F\times F$ object
by a $d\times d$ one, and combined with Proposition~\ref{prop:uniform-energy} it yields the
average linear error at every sparsity without a single Monte Carlo draw.

\subsection{Hyperparameters, complexity and cost}

The diagnostic has five hyperparameters: the readout rule, the threshold $\theta$ (default
$1/2$, the natural midpoint for unit-norm codes and Boolean states), the sparsity grid
$\mathcal S$, the trial budget $T$, and the block size $B$, which trades memory for the number
of regeneration passes. Only $T$ affects the statistical resolution;
Section~\ref{sec:res-e1} reports the threshold sensitivity of the conclusions.

Algorithm~\ref{alg:one} costs $O(Fd^2)$ time and $O(d^2)$ memory for the ratio $r$, by
Lemma~\ref{lem:suffstat}. The maximum-form ratio $r_{\max}$ is the one exception in the whole
method: no $d\times d$ reduction of a maximum is available to us, so obtaining it exactly costs
$O(F^2d)$ time and $O(F^2)$ memory, and it is computed only when asked for. Since the floor of
Theorem~\ref{thm:floor} is a statement about the mean square, nothing about the floor
comparison is lost by omitting it; we report $r_{\max}$ at the moderate widths where it is
affordable and drop it beyond. Algorithm~\ref{alg:two} costs $O(Fd^2+T\,F\,d\,s_{\max})$ time
and $O(Fd+d^2)$ memory. Algorithm~\ref{alg:ensemble} costs $O(T\,F\,d\,
s_{\max})$ time --- one pass of code regeneration and one $B\times d\times s_{\max}$ product
per block --- and $O(Bd+d^2)$ memory, independent of $F$. At $d=\EsixDmax$ and
$F=\EsixFmax$ the code would occupy $4$~GB in single precision; the streaming implementation
peaks near $B\,d\cdot4$ bytes. Measured wall-clock times are in Table~\ref{tab:cost}.

\subsection{Applicability, edge cases and failure modes}

\paragraph{Conditions of applicability.} The diagnostic assumes (i) $F>d$, otherwise the floor
is vacuous and Algorithm~\ref{alg:one} refuses to run; (ii) unit-norm code columns, or a
declared normalisation, since the threshold $\theta=1/2$ is calibrated to unit gain;
(iii) a sparse-state model that is exchangeable across coordinates, which is what makes the
closed form of step~3 valid.

\paragraph{Edge cases.} At $s=1$ the linear error is still non-zero while threshold recovery
succeeds whenever $\mu<\theta$; this is the smallest instance of the separation. When
$F/d\to1$ the floor tends to zero and the diagnosis becomes uninformative --- correctly so,
since a nearly complete code has no interference to speak of.

\paragraph{Failure modes.} Three are worth naming. (a) \emph{Degenerate gain}: if some
$|M_{ii}|$ is numerically zero the calibration is undefined and Algorithm~\ref{alg:one}
returns \textsc{Degenerate} rather than a number; this is not a corner case but the expected
outcome when a readout is fitted without the diagonal constraint
(Section~\ref{sec:res-e4}). (b) \emph{Uncovered feature}: a least-squares readout fitted on
states in which some feature never appears has an identically zero row, which triggers (a);
the implementation therefore enforces coverage of the fitting set. (c) \emph{Vacuous profile}:
if the coherence of the code exceeds $\theta$, $p_{\mathrm{rec}}(1)$ is already below $0.95$
and $s_{95}=0$; the certificate is then empty, and the diagnostic reports that rather than
extrapolating.

\subsection{Implementation and traceability}

\begin{table}[t]
\centering
\caption{Each step of the method and the function that implements it. Paths are relative to
the repository root.}
\label{tab:codemap}
\small
\begin{tabular}{lll}
\toprule
quantity & where defined & implementation \\
\midrule
$W(F,d)$ & Eq.~\eqref{eq:floor} & \code{lrtr/codes.py:welch\_floor} \\
$D^{-1}M$ & Cor.~\ref{cor:calibfree} & \code{lrtr/interface.py:unit\_diagonal} \\
$\widehat W,\widehat\mu,r,r_{\max}$ & Alg.~\ref{alg:one} & \code{lrtr/interface.py:crosstalk\_stats} \\
Algorithm~\ref{alg:one} & Sec.~\ref{sec:method} & \code{lrtr/diagnostic.py:interface\_floor\_diagnostic} \\
$\Sigma,v\mapsto\norm{A}_F^2,\norm{A\mathbf 1}^2$ & Lem.~\ref{lem:moments} & \code{lrtr/diagnostic.py:tied\_energy\_moments} \\
$\frac1F\E_S\norm{A\mathbf 1_S}^2$ & Prop.~\ref{prop:uniform-energy} & \code{lrtr/interface.py:linear\_energy\_uniform} \\
energy floor & Eq.~\eqref{eq:energyfloor} & \code{lrtr/interface.py:energy\_floor\_uniform} \\
threshold decoder $Q$ & Thm.~\ref{thm:coherence} & \code{lrtr/threshold.py:threshold\_decode} \\
fixed-code trial & Alg.~
ef{alg:two}, l.~5--7 & \code{lrtr/threshold.py:recovery\_trial\_fixed} \\
streaming trial & Alg.~
ef{alg:ensemble}, l.~6--10 & \code{lrtr/threshold.py:recovery\_trial\_streaming} \\
Algorithm~
ef{alg:two} & Sec.~
ef{sec:method} & \code{lrtr/diagnostic.py:fixed\_code\_separation\_profile} \\
Algorithm~
ef{alg:ensemble} & Sec.~
ef{sec:alg-ensemble} & \code{lrtr/diagnostic.py:random\_code\_scaling\_experiment} \\
$\operatorname{tr}(H\Sigma)-F$ & Lem.~
ef{lem:suffstat} & \code{lrtr/interface.py:crosstalk\_mean\_sq} \\
leverage form of $r$ & Prop.~
ef{prop:leverage} & \code{lrtr/interface.py:crosstalk\_mean\_sq\_pinv} \\
$s_{95}$ & Alg.~
ef{alg:two}, l.~9 & \code{lrtr/threshold.py:s95\_from\_curve} \\
\bottomrule
\end{tabular}
\end{table}

Table~\ref{tab:codemap} maps every quantity in Problem~\ref{prob:main} to the function that
computes it. The correctness of the implementation is checked by a test suite that verifies the
theorems numerically --- equality for tight frames, no violations for random interfaces, the
closed forms against Monte Carlo, and the streaming trial against an explicit dense evaluation
of the same code.

% ============================================================================
